% !TeX root = ../papers/nc-p3.tex
\section{大きな点ファイバーでの端支持と四文字ブロックの増大}
\label{sec:endpoint-plane-growth}
有限Frobenius窓の一次文字を $V_0,\ldots,V_3$ とし、$p\in V_3^*$ の三テンソル $w(p)$ の
端の支持を $V_0,V_2$ の部分空間とする。これらはそれぞれ始点・終点での隣接関係の評価像でもある。
辺の有限性、点列への延長可能性、または反対位置の対称性を仮定しない。
\begin{theorem}\label{thm:all-endpoint-ranks-two}
標準Koszul左右AS正則な四窓では、任意の非零一文字縮約の両端支持はともに次元2以上である。
従って全辺の全射影ファイバーは空、点、または射影直線であり、射影平面は現れない。
さらに二次元の端支持を与える点から、その支持平面への写像は各幾何学的ファイバーに高々一点を持つ。
\end{theorem}
\begin{lemma}\label{lem:rectangular-four-block}
$D\subset V_0$ の次元を $d$ とし、
\[
 k_D=\dim(R_0\cap(D\otimes V_1)),\qquad
 T=\{p\in V_3^*:w(p)\in D\otimes V_1\otimes V_2\},\quad l=\dim T
\]
とする。四位置を順に $V_0/D,V_1,V_2,T^*$ へ落とした商で、
長さ $4n$ の次元は $b^n$ であり
\[
 b\ \ge\ l(20-10d+4k_D-l)-k_D(10-4d+k_D).
\]
\end{lemma}
\begin{proof}
$R_3$ の像は零である。実際、$T$ での全評価が $D$ に入る。
従って四文字ごとの境界に関係はなく、一ブロックの次元を $b$ とすればよい。
第一の二文字の商は $4(4-d)-(6-k_D)=10-4d+k_D$ 次元である。
最初の三文字の周囲は $16(4-d)$ 次元、二つの挿入関係の次元は $4(6-k_D),6(4-d)$ である。
その交叉は元の三次重なりの射影を含み、これは $V_3^*/T$ に等しく $4-l$ 次元である。
よって三文字の商次元は $20-10d+4k_D-l$ 以上である。
$R_2$ の $V_2\otimes T^*$ への像は次元 $k_D$ 以下である。
実際、$w$ の第四文字を $T$ に制限すると、第一・第二文字側の全縮約は
$R_0\cap(D\otimes V_1)$ に入る。対応する隣接切断の他側の階数も同じである。
三文字商と $T^*$ の積から最後の関係像と最初の二文字商の積の次元上界を引けば結論を得る。
\end{proof}

\begin{proof}[一次元端支持の排除]
左端が $ku$ なら $w(p)=u\otimes M$ と書く。$M$ の階数を $r$、その左・右支持を $S,W$ とする。
$R_0$ は $u\otimes S$ を、$R_1$ は $M$ を含む。$r=1$ は連続する純関係により指数成長する。
$k=\dim(R_0\cap(ku\otimes V_1))$ とおく。$k=4$ なら $u$ が右単純のコピーを自由加群内に生成し、
ASの次数零Ext消滅に反する。$k=3$ なら $u\otimes S'\subset R_0$ となる三平面 $S'$ がある。
$R_1$ の $V_1/S'\otimes V_2$ への射影核は次元2以上であり、この核の各元に $u$ を左から付ければ
三次重なりとなる。従って補題の $d=1,k_D=3$ で $l\ge2$ である。
$2\le l\le4$ では $l(22-l)-27\ge13$ なので指数成長となる。

残りは $k=r=2$ としてよい。四位置を
$X=V_0/ku,V_1,Z=V_2/W,k_p$ に落とす。両端の $p$ に接する関係は零になる。
三文字 $X,V_1,Z$ の商の次元は、二つの関係像がそれぞれ4以下、5以下なので
$24-4\cdot2-3\cdot5=1$ に、それらの挿入の交叉次元を足した値以上である。
交叉は非零である。さもなくば全ての三次重なりが両端を落とした後で零となる。
すると任意の $0\ne z\in W^\perp\subset V_2^*$ で $w(z)$ の中央（位置0）の支持は $ku$ に入る。
非零の三テンソルで中央支持が一次元なら両側から同じ中央文字の純関係が得られ、
系\ref{cor:opposite-singular-rank-two}に反する。従って三文字商は二次元以上で、
$k_p$ により互いに独立なブロックを繰り返せば指数成長となる。
右端については語順を逆にする。
\end{proof}

\begin{proof}[二次元支持の一意性]
$d=2$ として $T\ne0$ とする。$w(p)$ の右端支持は次元2以上なので $k_D\ge2$ である。
一方、任意の $0\ne q\in D^\perp$ で $R_0(q)$ は次元2以上で、
この評価は $R_0/(R_0\cap(D\otimes V_1))$ を経由する。従って $k_D\le4$ である。
$2\le l\le4$, $2\le k_D\le4$ なら補題の下界
$l(4k_D-l)-k_D(2+k_D)$ は少なくとも4であり、標準成長に反する。
ゆえに $l=1$ である。
\end{proof}

\begin{corollary}\label{cor:endpoint-line-line-cut}
両端支持が二次元の $w(p)$ を持つ正則窓で、その支持を $D\subset V_0,W\subset V_2$ とすると、
$R_0\cap(D\otimes V_1)$ と $R_1\cap(V_1\otimes W)$ はともに二次元で、
ちょうど $w(p)$ の二端からの縮約像である。
両端を $D,W$ で割った全三次重なりの射影は高々一次元である。
\end{corollary}
\begin{proof}
二つの交叉次元を $k,h$ とすると各々2以上である。
$p$ を一点にし、他の位置を $V_0/D,V_1,V_2/W$ にしたとき、繰り返せる三文字商の次元は
$16-2(6-k)-2(6-h)+j=2(k+h)-8+j$ である。
$j$ は挿入した二つの関係の交叉次元で、全三次重なりの射影を含む。
この値は1以下なので $k=h=2$, $j\le1$ となる。
\end{proof}
